\documentclass[11pt,a4paper]{article}
\usepackage[utf8]{inputenc}
\usepackage[spanish]{babel}
\usepackage{amsmath, amssymb}
\usepackage{geometry}
\geometry{top=2.5cm, bottom=2.5cm, left=2.5cm, right=2.5cm}

\begin{document}

\section*{Piense, dialogue y comparta}

\begin{itemize}
    \item[\textbf{1.}] Usted está trabajando en una planta de fabricación. Una de las máquinas de la planta utiliza un resorte. Para un nuevo proceso, sería deseable que se utilizaran resortes con diferentes constantes elásticas durante diferentes partes del proceso. Su jefe le pide consejo sobre cómo cambiar un resorte por otro rápidamente para que todo el proceso pueda tener lugar en una cantidad razonable de tiempo. Usted sugiere que, en lugar de cambiar resortes continuamente, usted cambiaría la constante elástica de uno de los resortes largos sujetándolo en diferentes posiciones para definir un nuevo extremo fijo del resorte. Entonces el resorte efectivo consiste solamente en las bobinas que están más allá de la abrazadera. Usted diseña un sistema que consta de un resorte largo con $N$ bobinas y una constante elástica $k$. Usted diseña un sistema de abrazadera que aislará parte del resorte, dejando $N'$ bobinas más allá de la abrazadera fija. (a) Escriba una expresión para la constante elástica $k'$ del extremo libre del resorte en términos de $k, N$ y $N'$. (b) El extremo libre del resorte es sujetado y jalado hacia afuera una distancia $x$. En el proceso, la mano que sostiene el extremo del resorte hace un trabajo $W$ en el resorte. Ahora, se regresa el resorte a su estado relajado y luego se fija en su punto central. El extremo libre del resorte relajado con abrazadera, se sujeta y se jala hacia fuera la misma distancia $x$. ¿Cuánto trabajo hace la mano en el resorte en este caso? ¿$W$? ¿$2W$? ¿$4W$? ¿Otro valor?
    \vspace*{9cm}

    \item[\textbf{2.}] \textbf{ACTIVIDAD} En la tabla de datos, vemos las distancias de frenado mínimas $d$ en función de la rapidez inicial $v$ de un auto. Trabaje en su grupo para responder lo siguiente. (a) Si duplica la rapidez inicial, ¿se necesita el doble de distancia para detener el auto? (b) Suponga que la distancia de frenado es proporcional a la rapidez del auto elevada a cierta potencia: $d \propto v^n$. Usar técnicas gráficas para determinar $n$. (c) ¿Por qué la distancia de frenado depende del valor particular de $n$ que encontró en (b)?
    
    \begin{center}
    \begin{tabular}{cc}
    \hline
    \textbf{Rapidez (mi/h)} & \textbf{Distancia de frenado (ft)} \\
    \hline
    20 & 22.5 \\
    25 & 35.0 \\
    30 & 50.4 \\
    35 & 68.6 \\
    40 & 89.6 \\
    45 & 113.5 \\
    50 & 140.0 \\
    55 & 169.5 \\
    60 & 201.7 \\
    65 & 236.7 \\
    70 & 274.5 \\
    \hline
    \end{tabular}
    \end{center}
\end{itemize}

\newpage

\section*{Problemas}

\subsection*{SECCIÓN 7.2 Trabajo realizado por una fuerza constante}

\begin{itemize}
    \item[\textbf{1.}] Una cliente en un supermercado empuja un carro con una fuerza de $35.0 \text{ N}$ dirigida a un ángulo de $25.0^\circ$ bajo la horizontal. La fuerza es justo lo suficiente para equilibrar a varias fuerzas de fricción, así el carro se mueve con rapidez constante. (a) Encuentre el trabajo realizado por la cliente sobre el carro conforme ella avanza $50.0 \text{ m}$ a lo largo del pasillo. (b) La cliente camina al siguiente pasillo, empujando horizontalmente y manteniendo la misma rapidez que antes. Si la fuerza de fricción no cambia, ¿la fuerza aplicada de la cliente sería mayor, menor o la misma? (c) ¿Qué puede decirse sobre el trabajo realizado por la persona sobre el carro?
    \vspace*{6cm}

    \item[\textbf{2.}] El récord de levantamiento de botes, incluyendo el objeto y su tripulación de 10 miembros, corresponde a Sami Heinonen y Juha Räsänen de Suecia, en el año 2000. Ellos levantaron una masa total de $653.2 \text{ kg}$ aproximadamente a $4 \text{ pulgadas}$ del suelo, y lo hicieron 24 veces. Estime el trabajo total efectuado por los dos hombres en el bote en este levantamiento récord, ignore el trabajo negativo realizado por los hombres cuando descienden el bote al suelo.
    \vspace*{6cm}

    \item[\textbf{3.}] En 1990, Walter Arfeuille, de Bélgica, levantó un objeto de $281.5 \text{ kg}$ a través de una distancia de $17.1 \text{ cm}$ utilizando únicamente sus dientes. (a) ¿Cuánto trabajo realizó Arfeuille en este proceso, suponiendo que el objeto se levantó a rapidez constante? (b) ¿Qué fuerza total se ejerció en los dientes de Arfeuille durante este proceso?
    \vspace*{6cm}

    \item[\textbf{4.}] El hombre araña, cuya masa es $80.0 \text{ kg}$, está colgado en el extremo libre de una soga de $12.0 \text{ m}$, el otro extremo está fijo de la rama de un árbol arriba de él. Al flexionar repetidamente la cintura, hace que la cuerda se ponga en movimiento, y eventualmente la hace balancear lo suficiente para que pueda llegar a una repisa cuando la cuerda forma un ángulo de $60.0^\circ$ con la vertical. ¿Cuánto trabajo realizó la fuerza gravitacional en el hombre araña en esta maniobra?
    \vspace*{6cm}
\end{itemize}

\newpage

\subsection*{SECCIÓN 7.3 Producto escalar de dos vectores}

\begin{itemize}
    \item[\textbf{5.}] Para dos vectores cualesquiera $\vec{A}$ y $\vec{B}$, demuestre que $\vec{A} \cdot \vec{B} = A_x B_x + A_y B_y + A_z B_z$. \textit{Sugerencia:} Escriba $\vec{A}$ y $\vec{B}$ en forma de vectores unitarios y aplique las ecuaciones 7.4 y 7.5.
    \vspace*{6cm}

    \item[\textbf{6.}] El vector $\vec{A}$ tiene una magnitud de $5.00$ unidades y $\vec{B}$ tiene una magnitud de $9.00$ unidades. Los dos vectores forman un ángulo de $50.0^\circ$ entre sí. Determine $\vec{A} \cdot \vec{B}$.
    \vspace*{6cm}

    \item[\textbf{7.}] Encuentre el producto escalar de los vectores en la figura P7.7.
    \vspace*{6cm}

    \item[\textbf{8.}] Con la definición del producto escalar, encuentre los ángulos entre (a) $\vec{A} = 3\hat{i} - 2\hat{j}$ y $\vec{B} = 4\hat{i} - 4\hat{j}$, (b) $\vec{A} = -2\hat{i} + 4\hat{j}$ y $\vec{B} = 3\hat{i} - 4\hat{j} + 2\hat{k}$ y (c) $\vec{A} = \hat{i} - 2\hat{j} + 2\hat{k}$ y $\vec{B} = 3\hat{j} + 4\hat{k}$.
    \vspace*{6cm}
\end{itemize}

\newpage

\subsection*{SECCIÓN 7.4 Trabajo realizado por una fuerza variable}

\begin{itemize}
    \item[\textbf{9.}] Una partícula se somete a una fuerza $F_x$ que varía con la posición como se muestra en la figura P7.9. Encuentre el trabajo realizado por la fuerza en la partícula mientras se mueve (a) de $x = 0$ a $x = 5.00 \text{ m}$, (b) de $x = 5.00 \text{ m}$ a $x = 10.0 \text{ m}$ y (c) de $x = 10.0 \text{ m}$ a $x = 15.0 \text{ m}$. (d) ¿Cuál es el trabajo total efectuado por la fuerza a través de la distancia $x = 0$ a $x = 15.0 \text{ m}$?
    \vspace*{6cm}

    \item[\textbf{10.}] En un sistema de control, un acelerómetro consiste en un objeto de $4.70 \text{ g}$ que se desliza sobre un riel horizontal calibrado. Un resorte de masa pequeña une al objeto a una pestaña en un extremo del riel. La grasa en el riel hace despreciable la fricción estática, pero amortigua rápidamente las vibraciones del objeto deslizante. Cuando se somete a una aceleración fija de $0.800g$, el objeto debería estar en una posición $0.500 \text{ cm}$ de su posición de equilibrio. Encuentre la constante elástica del resorte requerida para una correcta calibración.
    \vspace*{6cm}

    \item[\textbf{11.}] Cuando un objeto de $4.00 \text{ kg}$ cuelga verticalmente en cierto resorte ligero descrito por la ley de Hooke, el resorte se estira $2.50 \text{ cm}$. Si se quita el objeto de $4.00 \text{ kg}$, (a) ¿cuánto se estirará el resorte si se le cuelga un objeto de $1.50 \text{ kg}$? (b) ¿Cuánto trabajo debe efectuar un agente externo para estirar el mismo resorte $4.00 \text{ cm}$ desde su posición sin estirar?
    \vspace*{6cm}

    \item[\textbf{12.}] Exprese, en unidades fundamentales SI, las unidades de la constante elástica de un resorte.
    \vspace*{5cm}
\end{itemize}

\newpage

\begin{itemize}
    \item[\textbf{13.}] El dispensador de charolas en su cafetería se ha roto y no se puede reparar. El encargado sabe que usted es bueno diseñando cosas y le pide que le ayude a construir un nuevo dispensador con partes que tiene en su mesa de trabajo. El dispensador de charolas sostiene una pila de estas sobre un anaquel que se apoya en cuatro resortes en uno en cada esquina del anaquel. Cada charola es rectangular, de $45.3 \text{ cm}$ por $35.6 \text{ cm}$, $0.450 \text{ cm}$ de grosor y $580 \text{ g}$ de masa. El encargado le pide que diseñe un nuevo dispensador de cuatro resortes de manera que cuando se saca una charola del dispensador este empuja hacia arriba a la pila restante para que la charola superior esté en la misma posición que la de la charola que acaba de quitar. Él tiene una gran variedad de resortes que puede utilizar para construir el dispensador. ¿Qué resortes debe usar?
    \vspace*{7cm}

    \item[\textbf{14.}] Un resorte ligero, con constante elástica de $3.85 \text{ N/m}$, se comprime $8.00 \text{ cm}$ mientras se mantiene entre un bloque de $0.250 \text{ kg}$ a la izquierda y un bloque de $0.500 \text{ kg}$ a la derecha, ambos en reposo sobre una superficie horizontal. El resorte ejerce una fuerza en cada bloque, y tiende a separarlos. Los bloques se sueltan simultáneamente desde el reposo. Encuentre la aceleración con la que cada bloque comienza a moverse, dado que el coeficiente de fricción cinética entre cada bloque y la superficie es (a) $0$, (b) $0.100$ y (c) $0.462$.
    \vspace*{7cm}

    \item[\textbf{15.}] Una partícula pequeña de masa $m$ se jala hacia lo alto de un medio cilindro sin fricción (de radio $R$) mediante una cuerda que pasa sobre lo alto del cilindro, como se ilustra en la figura P7.15. (a) Si supone que la partícula se mueve con rapidez constante, demuestre que $F = mg \cos\theta$. \textit{Nota:} Si la partícula se mueve con rapidez constante, la componente de su aceleración tangente al cilindro debe ser cero en todo momento. (b) Mediante integración directa de $W = \int \vec{F} \cdot d\vec{r}$, encuentre el trabajo realizado al mover la partícula con rapidez constante desde el fondo hasta lo alto del medio cilindro.
    \vspace*{7cm}
\end{itemize}

\newpage

\begin{itemize}
    \item[\textbf{16.}] La fuerza que actúa en una partícula es $F_x = (8x - 16)$, donde $F$ está en newtons y $x$ está en metros. (a) Grafique esta fuerza en términos de $x$ desde $x = 0$ hasta $x = 3.00 \text{ m}$. (b) A partir de su gráfica, encuentre el trabajo neto realizado por esta fuerza en la partícula conforme se traslada de $x = 0$ a $x = 3.00 \text{ m}$.
    \vspace*{6cm}

    \item[\textbf{17.}] Cuando se cuelgan diferentes cargas de un resorte, el resorte se estira a diferentes longitudes, como se muestra en la siguiente tabla. (a) Elabore una gráfica de la fuerza aplicada en términos de la extensión del resorte. (b) Mediante ajuste por mínimos cuadrados, determine la línea recta que ajusta mejor los datos. (c) Para completar el inciso (b), ¿quiere emplear todos los datos o debería ignorar algunos de ellos? Explique. (d) A partir de la pendiente de la línea recta de mejor ajuste, encuentre la constante elástica $k$. (e) Si el resorte se extiende a $105 \text{ mm}$, ¿qué fuerza ejerce sobre el objeto suspendido?
    
    \begin{center}
    \begin{tabular}{lccccccccccc}
    \hline
    $F \text{ (N)}$ & 2.0 & 4.0 & 6.0 & 8.0 & 10 & 12 & 14 & 16 & 18 & 20 & 22 \\
    $L \text{ (mm)}$ & 15 & 32 & 49 & 64 & 79 & 98 & 112 & 126 & 149 & 175 & 190 \\
    \hline
    \end{tabular}
    \end{center}
    \vspace*{7cm}

    \item[\textbf{18.}] Se dispara una bola de $100 \text{ g}$ de un rifle que tiene un cañón de $0.600 \text{ m}$ de largo. Elija el origen como la ubicación donde la bala comienza a moverse. Entonces la fuerza (en newtons) que ejercen en la bala los gases en expansión es $15\,000 + 10\,000x - 25\,000x^2$, donde $x$ está en metros. (a) Determine el trabajo efectuado por el gas en la bala conforme esta recorre la longitud del cañón. (b) \textbf{¿Qué pasaría si?} Si el cañón mide $1.00 \text{ m}$ de largo, ¿cuánto trabajo se realiza? y (c) ¿Cómo se compara este valor con el trabajo calculado en el inciso (a)?
    \vspace*{6cm}

    \item[\textbf{19.}] Una fuerza $\vec{F} = (4x\hat{i} + 3y\hat{j})$, donde $\vec{F}$ está en newtons y $x$ y $y$ están en metros, actúa en un objeto conforme este se mueve en la dirección $x$ desde el origen hasta $x = 5.00 \text{ m}$. Encuentre el trabajo $W = \int \vec{F} \cdot d\vec{r}$ efectuado por la fuerza en el objeto. (b) \textbf{¿Qué pasaría si?} Encuentre el trabajo $W = \int \vec{F} \cdot d\vec{r}$ hecho por la fuerza en el objeto si se mueve del origen a $(5.0 \text{ m}, 5.00 \text{ m})$ a lo largo de una trayectoria recta que hace un ángulo de $45.0^\circ$ con el eje $x$ positivo. ¿Es el trabajo realizado por esta fuerza dependiente de la trayectoria tomada entre los puntos iniciales y finales?
    \vspace*{6cm}

    \item[\textbf{20.}] \textbf{Problema de repaso.} La gráfica en la figura P7.20 especifica una relación entre $u$ y $v$. (a) Evalúe $\int_0^v u\,dv$. (b) Obtenga $\int_0^a u\,dv$. (c) Encuentre $\int_a^b v\,du$.
    \vspace*{6cm}
\end{itemize}

\newpage

\subsection*{SECCIÓN 7.5 Energía cinética y teorema trabajo-energía cinética}

\begin{itemize}
    \item[\textbf{21.}] Una partícula de $0.600 \text{ kg}$ tiene una rapidez de $2.00 \text{ m/s}$ en el punto \textcircled{A} y energía cinética de $7.50 \text{ J}$ en el punto \textcircled{B}. ¿Cuál es (a) su energía cinética en \textcircled{A}, (b) su rapidez en \textcircled{B} y (c) el trabajo neto realizado en la partícula por las fuerzas externas cuando se mueve de \textcircled{A} a \textcircled{B}?
    \vspace*{6cm}

    \item[\textbf{22.}] Una partícula de $4.00 \text{ kg}$ se somete a una fuerza neta que varía con la posición, como se muestra en la figura P7.15. La partícula parte del reposo en $x = 0$. ¿Cuál es su rapidez en (a) $x = 5.00 \text{ m}$, (b) $x = 10.0 \text{ m}$ y (c) $x = 15.0 \text{ m}$?
    \vspace*{6cm}

    \item[\textbf{23.}] Un martinete de $2\,100 \text{ kg}$ se usa para enterrar una viga de acero en la tierra. El martinete cae $5.00 \text{ m}$ antes de quedar en contacto con la parte superior de la viga, y clava la viga $12.0 \text{ cm}$ más en el suelo mientras llega al reposo. Aplicando consideraciones de energía, calcule la fuerza promedio que la viga ejerce en el martinete mientras éste logra el reposo.
    \vspace*{6cm}

    \item[\textbf{24.}] \textbf{Problema de repaso.} En un microscopio electrónico, hay un cañón de electrones que contiene dos placas metálicas con carga, separadas $2.80 \text{ cm}$. Una fuerza eléctrica acelera cada electrón en el haz desde el reposo hasta $9.60\%$ de la rapidez de la luz a través de esta distancia. (a) Determine la energía cinética del electrón mientras abandona el cañón de electrones. Los electrones llevan esta energía a un material fosforescente en la superficie interior de la pantalla del televisor y lo hacen brillar. Para un electrón que pasa entre las placas en el cañón de electrones, determine, (b) la magnitud de la fuerza eléctrica constante que actúa en el electrón, (c) la aceleración del electrón y (d) el tiempo de vuelo entre las placas.
    \vspace*{7cm}
\end{itemize}

\newpage

\begin{itemize}
    \item[\textbf{25.}] \textbf{Problema de repaso.} Se puede considerar al teorema trabajo-energía cinética como una segunda teoría de movimiento, paralela a las leyes de Newton, en cuanto que describe cómo las influencias externas afectan el movimiento de un objeto. En este problema, resuelva los incisos (a), (b) y (c) por separado de los incisos (d) y (e), de modo que pueda comparar las predicciones de las dos teorías. En un cañón de rifle de longitud $72.0 \text{ cm}$, una bala de $15.0 \text{ g}$ acelera desde el reposo hasta una rapidez de $780 \text{ m/s}$. (a) Encuentre la energía cinética de la bala cuando abandona el cañón. (b) Use el teorema trabajo-energía cinética para obtener el trabajo neto realizado sobre la bala. (c) Utilice su resultado del inciso (b) para encontrar la magnitud de la fuerza neta promedio que actuó en la bala mientras esta estaba en el cañón. (d) Ahora modele a la bala como una partícula bajo aceleración constante. Obtenga la aceleración constante de una bala que inicia desde el reposo y gana una rapidez de $780 \text{ m/s}$ en una distancia de $72.0 \text{ cm}$. (e) Modele la bala como una partícula bajo una fuerza neta, para encontrar la fuerza neta que actuó en ella durante su aceleración. (f) ¿Qué conclusión obtiene al comparar los resultados de los incisos (c) y (e)?
    \vspace*{8cm}

    \item[\textbf{26.}] Usted está acostado en su dormitorio, descansando después de hacer su tarea de física. A medida que mira fijamente el techo, le llega la idea de un nuevo juego. Toma un dardo con una nariz pegajosa y una masa de $19.0 \text{ g}$. También toma un resorte que tiene en su escritorio de algún proyecto anterior. Pinta un objetivo en el techo. Su nuevo juego es colocar el resorte verticalmente en el suelo, colocar el dardo de la nariz pegajosa hacia arriba en el resorte, y empujar el resorte hacia abajo hasta que el resorte esté comprimido, como se muestra a la derecha en la figura P7.26. Luego suelta el resorte, disparando el dardo hacia el blanco en su techo, donde su nariz pegajosa hará que cuelgue del techo. El resorte tiene una longitud sin comprimir de extremo a extremo de $5.00 \text{ cm}$, como se muestra a la izquierda en la figura P7.26, se puede comprimir a una longitud de extremo a extremo de $1.00 \text{ cm}$ cuando el resorte está muy comprimido. Antes de probar el juego, usted sostiene el extremo superior del resorte en una mano y cuelga un paquete de 10 dardos idénticos del extremo inferior del resorte. El resorte se extiende $1.00 \text{ cm}$ debido al peso de los dardos. Está tan entusiasmado con el nuevo juego que, antes de hacer una prueba del juego, corre para reunirse con sus amigos para mostrarlo. Cuando sus amigos están en su habitación y les muestra el primer disparo de su nuevo juego, ¿por qué se avergüenza?
    \vspace*{7cm}

    \item[\textbf{27.}] \textbf{Problema de repaso.} Un objeto de $5.75 \text{ kg}$ pasa por el origen en el tiempo $t = 0$ tal que la componente $x$ de su velocidad es $5.00 \text{ m/s}$ y su componente $y$ de velocidad es $-3.00 \text{ m/s}$. (a) ¿Cuál es la energía cinética del objeto en este instante? (b) Al tiempo $t = 2.00 \text{ s}$, la partícula se localiza en $x = 8.50 \text{ m}$ y $y = 5.00 \text{ m}$. ¿Qué fuerza constante actuó sobre el objeto durante este intervalo de tiempo? (c) ¿Cuál es la rapidez de la partícula en $t = 2.00 \text{ s}$?
    \vspace*{7cm}

    \item[\textbf{28.}] \textbf{Problema de repaso.} Una bala de $7.80 \text{ g}$ con una rapidez de $575 \text{ m/s}$ golpea la mano de un superhéroe, haciendo que su mano se mueva $5.50 \text{ cm}$ en la dirección de la velocidad de la bala antes de parar. (a) Use consideraciones de trabajo y energía para encontrar la fuerza promedio que detuvo la bala. (b) Suponga que la fuerza es constante, luego determine cuánto tiempo transcurrió entre el momento que la bala golpeó la mano hasta el instante en que se detiene.
    \vspace*{6cm}
\end{itemize}

\newpage

\subsection*{SECCIÓN 7.6 Energía potencial de un sistema}

\begin{itemize}
    \item[\textbf{29.}] Una piedra de $0.20 \text{ kg}$ se mantiene a $1.3 \text{ m}$ sobre el borde de un pozo de agua y entonces se le deja caer en él. El pozo tiene una profundidad de $5.0 \text{ m}$. Respecto a la configuración con la piedra en el borde del pozo, ¿cuál es la energía potencial gravitacional del sistema piedra-Tierra (a) antes de liberar la piedra y (b) cuando esta llega al fondo del pozo? (c) ¿Cuál es el cambio en energía potencial gravitacional del sistema desde que se suelta la piedra hasta que alcanza el fondo del pozo?
    \vspace*{6cm}

    \item[\textbf{30.}] Un carro de montaña rusa, de $1\,000 \text{ kg}$, inicialmente está en lo alto de un bucle, en el punto \textcircled{A}. Luego se mueve $135 \text{ pies}$ a un ángulo de $40.0^\circ$ bajo la horizontal, hacia un punto inferior \textcircled{B}. (a) Elija el carro en el punto \textcircled{B} como la configuración cero para la energía potencial gravitacional del sistema montaña rusa-Tierra. Determine la energía potencial del sistema cuando el carro está en los puntos \textcircled{A} y \textcircled{B}, y el cambio en energía potencial conforme se mueve el carro entre estos puntos. (b) Repita el inciso (a), pero haga la configuración cero con el carro en el punto \textcircled{A}.
    \vspace*{6cm}
\end{itemize}

\subsection*{SECCIÓN 7.7 Fuerzas conservativas y no conservativas}

\begin{itemize}
    \item[\textbf{31.}] Una partícula de $4.00 \text{ kg}$ se mueve desde el origen a la posición \textcircled{C}, que tiene coordenadas $x = 5.00 \text{ m}$ y $y = 5.00 \text{ m}$ (figura P7.31). Una fuerza en la partícula es la fuerza gravitacional que actúa en la dirección $y$ negativa. Con la ecuación 7.3, calcule el trabajo realizado por la fuerza gravitacional en la partícula conforme va de $O$ a \textcircled{C} a lo largo de (a) la trayectoria violeta, (b) la trayectoria roja y (c) la trayectoria azul. (d) Sus resultados deberían ser idénticos, ¿por qué?
    \vspace*{6cm}

    \item[\textbf{32.}] (a) Suponga que una fuerza constante actúa en un objeto. La fuerza no varía con el tiempo o con la posición o la velocidad del objeto. Comience con la definición general del trabajo realizado por una fuerza $W = \int \vec{F} \cdot d\vec{r}$ y demuestre que la fuerza es conservativa. (b) Como caso especial, suponga que la fuerza $\vec{F} = (3\hat{i} + 4\hat{j})$ actúa en una partícula que se mueve de $O$ a \textcircled{C} en la figura P7.31. Calcule el trabajo efectuado por $\vec{F}$ en la partícula conforme se mueve a lo largo de cada una de las tres trayectorias mostradas en la figura y demuestre que el trabajo realizado tiene el mismo valor para las tres trayectorias. (c) \textbf{¿Qué pasaría si?} ¿Es el trabajo hecho también idéntico a lo largo de las tres trayectorias de la fuerza $\vec{F} = (4x\hat{i} + 3y\hat{j})$, donde $\vec{F}$ está en Newtons y $x$ y $y$ están en metros, del problema 19? (d) \textbf{¿Qué pasaría si?} Suponga que la fuerza está dada por $\vec{F} = (y\hat{i} - x\hat{j})$, donde $\vec{F}$ está en newtons y $x$ y $y$ está en metros. ¿Es el trabajo hecho idéntico a lo largo de las tres trayectorias para esta fuerza?
    \vspace*{8cm}
\end{itemize}

\newpage

\begin{itemize}
    \item[\textbf{33.}] Una fuerza que actúa en una partícula moviéndose en el plano $xy$ está dada por $\vec{F} = (2y\hat{i} + x^2\hat{j})$, donde $\vec{F}$ está en newtons y $x$ y $y$ están en metros. La partícula se mueve desde el origen hasta la posición final con coordenadas $x = 5.00 \text{ m}$ y $y = 5.00 \text{ m}$, como se muestra en la figura P7.31. Calcule el trabajo efectuado por $\vec{F}$ sobre la partícula cuando esta se desplaza a lo largo de (a) la trayectoria violeta, (b) la trayectoria roja y (c) la trayectoria azul (véase su respectivo QR). (d) ¿$\vec{F}$ es conservativa o no conservativa? (e) Explique su respuesta al inciso (d).
    \vspace*{8cm}
\end{itemize}

\subsection*{SECCIÓN 7.8 Relación entre fuerzas conservativas y energía potencial}

\begin{itemize}
    \item[\textbf{34.}] ¿Por qué no es posible la siguiente situación? Un bibliotecario levanta un libro desde el suelo hasta un estante, haciendo $20.0 \text{ J}$ de trabajo en este proceso. Cuando él se voltea, el libro cae del estante de regreso al suelo. La fuerza gravitacional de la Tierra sobre el libro efectúa $20.0 \text{ J}$ de trabajo sobre él mientras cae. Como el trabajo realizado fue $20.0 \text{ J} + 20.0 \text{ J} = 40.0 \text{ J}$, el libro golpea el suelo con $40.0 \text{ J}$ de energía cinética.
    \vspace*{6cm}

    \item[\textbf{35.}] Una sola fuerza conservativa actúa en una partícula de $5.00 \text{ kg}$ dentro de un sistema debido a su interacción con el resto del sistema. La ecuación $F_x = 2x + 4$ describe la fuerza, donde $F_x$ está en newtons y $x$ está en metros. Conforme la partícula se desplaza a lo largo del eje $x$ desde $x = 1.00 \text{ m}$ hasta $x = 5.00 \text{ m}$, calcule (a) el trabajo realizado por esta fuerza en la partícula, (b) el cambio en la energía potencial del sistema y (c) la energía cinética que tiene la partícula en $x = 5.00 \text{ m}$ si su rapidez es $3.00 \text{ m/s}$ en $x = 1.00 \text{ m}$.
    \vspace*{6cm}

    \item[\textbf{36.}] Una función de energía potencial para un sistema en el cual una fuerza en dos dimensiones actúa de la forma $U = 3x^3y - 7x$. Encuentre la fuerza que actúa en el punto $(x, y)$.
    \vspace*{6cm}

    \item[\textbf{37.}] La energía potencial de un sistema de dos partículas separadas por una distancia $r$ está dada por $U(r) = A/r$, donde $A$ es una constante. Encuentre la fuerza radial $\vec{F}_r$ que cada partícula ejerce sobre la otra.
    \vspace*{6cm}
\end{itemize}

\newpage

\subsection*{SECCIÓN 7.9 Diagramas de energía y equilibrio de un sistema}

\begin{itemize}
    \item[\textbf{38.}] Para la curva de energía potencial que se muestra en la figura P7.38, (a) determine si la fuerza $F_x$ es positiva, negativa o cero en los cinco puntos señalados. (b) Indique los puntos de equilibrio estable, inestable y neutro. (c) Trace la gráfica de la curva de $F_x$ en función de $x$ desde $x = 0$ hasta $x = 9.5 \text{ m}$.
    \vspace*{7cm}

    \item[\textbf{39.}] Un cono circular recto se puede equilibrar, teóricamente, sobre una superficie horizontal en tres diferentes formas. Trace estas tres configuraciones de equilibrio e identifíquelas como posiciones de equilibrio estable, inestable o neutro.
    \vspace*{6cm}
\end{itemize}

\section*{PROBLEMAS ADICIONALES}

\begin{itemize}
    \item[\textbf{40.}] La función de energía potencial de un sistema de partículas está dada por $U(x) = -x^3 + 2x^2 + 3x$, donde $x$ es la posición de una partícula en el sistema. (a) Determine la fuerza $F_x$ en la partícula como una función de $x$. (b) ¿Para qué valores de $x$ es la fuerza igual a cero? (c) Grafique $U(x)$ en función de $x$ y $F_x$ en términos de $x$ y señale los puntos de equilibrio estable y de equilibrio inestable.
    \vspace*{7cm}

    \item[\textbf{41.}] Usted tiene una nueva pasantía, donde está ayudando a diseñar un nuevo patio de carga para la estación de trenes de su ciudad. Habrá un número de vías de carga donde se pueden almacenar vagones hasta que se necesiten. Para evitar que los vagones se deslicen en las vías, se ha diseñado una combinación de dos resortes en espiral como se ilustra en la figura P7.41. Cuando un vagón se mueve a la derecha en la figura y pega en los resortes, que ejercen una fuerza hacia la izquierda sobre el vagón para frenarlo. Ambos resortes se describen con la ley de Hooke y tienen constantes elásticas $k_1 = 1\,600 \text{ N/m}$ y $k_2 = 3\,400 \text{ N/m}$. Después de que el primer resorte se comprime una distancia de $30.0 \text{ cm}$, el segundo resorte actúa con el primero para aumentar la fuerza hacia la izquierda en el vagón en la figura P7.41. Cuando el resorte con la constante elástica $k_2$ se comprime $50.0 \text{ cm}$, las espirales de alambre de ambos resortes se presionan entre sí, por lo que los resortes ya no se pueden comprimir. Un vagón típico en la vía de carga tiene una masa de $6\,000 \text{ kg}$. Cuando usted presenta su diseño a su supervisor, le pide la rapidez máxima que un vagón puede tener y ser detenido por su dispositivo.
    \vspace*{7cm}
\end{itemize}

\newpage

\begin{itemize}
    \item[\textbf{42.}] Cuando un objeto se desplaza una cantidad $x$ de su equilibrio estable, entonces una fuerza restauradora actúa en él, intentando retornar al objeto a su posición de equilibrio. La magnitud de la fuerza de restitución puede ser una complicada función de $x$. En tales casos, en general se puede suponer que la función de fuerza $F(x)$ es una serie de potencias en $x$ como $F(x) = -(k_1x + k_2x^2 + k_3x^3 + \dots)$. Aquí el primer término es justamente la ley de Hooke, la cual describe la fuerza ejercida por un resorte simple para pequeños desplazamientos. Para pequeñas excursiones fuera del equilibrio, en general ignoramos los términos de orden superior, pero en algunos casos puede ser deseable mantener el segundo término. Si la fuerza de restitución se modela como $F(x) = -(k_1x + k_2x^2)$, ¿cuánto trabajo se realiza en un objeto al desplazarlo de $x = 0$ a $x = x_{\text{máx}}$ mediante una fuerza aplicada $-F$?
    \vspace*{6cm}

    \item[\textbf{43.}] Una partícula se mueve a lo largo del eje $x$ desde $x = 12.8 \text{ m}$ hasta $x = 23.7 \text{ m}$ sometido a la influencia de una fuerza $F = \frac{375}{x^2 + 3.75x}$ donde $F$ está en newtons y $x$ en metros. Con el uso de integración numérica, determine el trabajo efectuado por esta fuerza en la partícula durante este desplazamiento. Su resultado debe ser preciso dentro de $2\%$.
    \vspace*{6cm}

    \item[\textbf{44.}] ¿Por qué no es posible la siguiente situación? En un nuevo casino, se introduce una máquina de \textit{pinball} gigante. El casino advierte que su tamaño es comparable a la altura de un jugador de basketball. El lanzador de bola en la máquina envía bolas de metal por un lado de la máquina para jugar. El resorte en el lanzador (figura P7.44) tiene una constante elástica de $1.20 \text{ N/cm}$. La superficie sobre la que se mueve la bola está inclinada $\theta = 10.0^\circ$ respecto de la horizontal. El resorte inicialmente se comprime su distancia máxima $d = 5.00 \text{ cm}$. Una bola de masa $100 \text{ g}$ es lanzada para entrar en juego al soltar el émbolo. En el casino los visitantes encuentran muy emocionante este juego con la gigantesca máquina.
    \vspace*{6cm}

    \item[\textbf{45.}] \textbf{Problema de repaso.} Dos fuerzas constantes actúan en un objeto de $5.00 \text{ kg}$ que se mueve en el plano $xy$, como se muestra en la figura P7.45. La fuerza $\vec{F}_1$ es de $25.0 \text{ N}$ a $35.0^\circ$ y la fuerza $\vec{F}_2$ es de $42.0 \text{ N}$ a $150^\circ$. En el tiempo $t = 0$, el objeto está en el origen y tiene velocidad $(4.00\hat{i} + 2.50\hat{j})$. (a) Exprese las dos fuerzas en notación de vector unitario. Use notación de vectores unitarios para sus otras respuestas. (b) Encuentre la fuerza total que se ejerce en el objeto. (c) Encuentre la aceleración del objeto. Ahora, considere el instante $t = 3.00 \text{ s}$, encuentre (d) la velocidad del objeto, (e) su posición, (f) su energía cinética a partir de $\frac{1}{2}mv_f^2$ y (g) su energía cinética a partir de $\frac{1}{2}mv_i^2 + \sum \vec{F} \cdot \Delta\vec{r}$. (h) ¿Qué conclusión puede sacar al comparar las respuestas a los incisos (f) y (g)?
    \vspace*{8cm}
\end{itemize}

\newpage

\begin{itemize}
    \item[\textbf{46.}] (a) Tome $U = 5$ para un sistema con una partícula en la posición $x = 0$ y calcule la energía potencial del sistema como una función de la posición $x$ de la partícula. La fuerza sobre la partícula está dada por $(8e^{-2x})\hat{i}$. (b) Explique si la fuerza es conservativa o no conservativa y cómo puede decirlo.
    \vspace*{6cm}

    \item[\textbf{47.}] Un plano inclinado de ángulo $\theta = 20.0^\circ$ tiene un resorte de constante elástica $k = 500 \text{ N/m}$ anclado en el fondo tal que el resorte queda paralelo a la superficie, como se muestra en la figura P7.47. Un bloque de masa $m = 2.50 \text{ kg}$ se coloca en el plano a una distancia $d = 0.300 \text{ m}$ del resorte. Desde esta posición, el bloque es proyectado hacia abajo dirigido al resorte con rapidez $v = 0.750 \text{ m/s}$. ¿Cuánto se comprime el resorte cuando el bloque está momentáneamente en reposo?
    \vspace*{6cm}

    \item[\textbf{48.}] Un plano inclinado de ángulo $\theta$ tiene un resorte de constante elástica $k$ anclado en el fondo tal que el resorte es paralelo a la superficie. Un bloque de masa $m$ se coloca sobre el plano a una distancia $d$ del resorte. Desde esta posición, se libera el bloque y se proyecta hacia el resorte con rapidez $v$, como se muestra en la figura P7.47. ¿Cuánto se comprime el resorte cuando el bloque está momentáneamente en reposo?
    \vspace*{6cm}

    \item[\textbf{49.}] Durante las vacaciones de Navidad, está haciendo algo de dinero extra para comprar regalos trabajando en una fábrica, ayudando a mover cajas. En un momento determinado, usted encuentra que todas las carretillas, plataformas y carros están en uso, por lo que debe mover una caja a través de la habitación una distancia en línea recta de $35.0 \text{ m}$ sin la ayuda de estos dispositivos. Observe que la caja tiene una cuerda atada a la mitad de una de sus caras verticales. Decide mover la caja jalando de la soga. La caja tiene una masa de $130 \text{ kilogramos}$, y el coeficiente de fricción cinética entre la caja y el piso de concreto es $0.350$. (a) Determine el ángulo respecto a la horizontal con el que debe jalar hacia arriba la cuerda para que pueda mover la caja la distancia deseada con la fuerza de magnitud \textit{más pequeña}. (b) Con este ángulo jale la cuerda, ¿cuánto trabajo haces al arrastrar la caja en la distancia deseada?
    \vspace*{7cm}
\end{itemize}

\newpage

\subsection*{PROBLEMA DE DESAFÍO}

\begin{itemize}
    \item[\textbf{50.}] Una partícula de masa $m = 1.18 \text{ kg}$ se une entre dos resortes idénticos sobre una mesa horizontal sin fricción. Ambos resortes tienen constante elástica $k$ e inicialmente no están estirados, y la partícula está en $x = 0$. (a) La partícula se jala una distancia $x$ a lo largo de una dirección perpendicular a la configuración inicial de los resortes, como se muestra en la figura P7.50. Demuestre que la fuerza ejercida en la partícula es
    $$ \vec{F} = -2kx \left( 1 - \frac{L}{\sqrt{x^2 + L^2}} \right)\hat{i} $$
    (b) Demuestre que la energía potencial del sistema es
    $$ U(x) = kx^2 + 2kL(L - \sqrt{x^2 + L^2}) $$
    (c) Elabore una gráfica de $U(x)$ en términos de $x$ e identifique todos los puntos de equilibrio. Suponga $L = 1.20 \text{ m}$ y $k = 40.0 \text{ N/m}$. (d) Si la partícula se jala $0.500 \text{ m}$ hacia la derecha y después se libera, ¿cuál es su rapidez al llegar a $x = 0$?
    \vspace*{9cm}
\end{itemize}

\end{document}
